% ======================================================================
% SECTION 4: FINDINGS - DRAFT REORGANISATION
% ======================================================================
% Working draft for pipeline-stage structure
% Source: /home/shawn/Code/LLM-History-Paper/AbsenceJudgement.tex
% Plan: /home/shawn/Code/theseus-ship/llm-absence-judgement/planning/section4-reorganisation-plan.md
%
% Batch 1: 4.1 Overview (no changes) + 4.2 Ideation (new)
% ======================================================================

\section{Findings: AI Research Tool Case Studies}

% ----------------------------------------------------------------------
% 4.1 OVERVIEW AND SCOPE - NO CHANGES
% ----------------------------------------------------------------------
% Keep existing content from lines 348-373
% FLAG: Table 1 (tab:tool-usage) needs update to reflect actual model usage per stage

\subsection{Overview and Scope}

% [KEEP EXISTING CONTENT - no changes required]
% Table 1 flagged for later update

% ----------------------------------------------------------------------
% 4.2 IDEATION AND RESEARCH DESIGN - NEW SECTION
% ----------------------------------------------------------------------
% Extracted from: Claude Ideation (lines 433-436), Gemini Ideation (lines 478-482),
%                 Appendix (line 745), chat archives (methodology, Google section 4)
% NOTE: OpenAI Synthesis content (lines 404-406) moved to 4.3 Literature Discovery

\subsection{Ideation and Research Design}

We evaluated GPT-4.5, Claude 3.7 Sonnet, and Gemini 2.5 Pro for ideation tasks including research question generation, framework development, and outlining. Using the AI-facilitated self-interview approach described in Section~\ref{sec:method}, we found that model performance varied, and that different models ``asked different questions''. While this category of use cannot be measured as objectively as literature discovery or data acquisition, we identified informative strengths and weaknesses across models.

\subsubsection{OpenAI}

GPT-4.5 was used for initial ideation and outlining, occasionally producing unexpected insights that other models missed. The o3 model's approach to research question generation was formulaic, and its outputs adopted language unsuitable for academic writing---phrases like ``accelerate clarity'' that one author characterised as ``corporate marketing''. 

\subsubsection{Anthropic}

Claude 3.7 Sonnet demonstrated better performance in generating research questions. When tasked with suggesting frameworks for analysis, it produced more nuanced categories and showed awareness of interdisciplinary considerations. Claude better matched the stylistic register required for scholarly work. When prompted to ask focused questions, Claude retained focus and sought depth in the context of the topic rather than generating lengthy but shallow question sets. These traits made Sonnet the preferred tool for iterative development through structured questioning.

\subsubsection{Google}

Gemini 2.5 Pro's questions proved superior to o3's formulaic approach but lacked both the consistency of Claude 3.7 Sonnet and the occasional insights of GPT-4.5. In one instance, Gemini confused human and AI identities in its opening message, referring to ``my autoethnographic findings'' as if it were the researcher. When asked to help develop Section 4's outline after being provided with 5,500 words explicitly discussing the paper's central argument, Gemini posed the question: ``what is the overarching argument you intend to make?'' Combined with Gemini's tendency to conflate multiple complex tasks into single operations (discussed further in Section~\ref{sec:lit-discovery}), these limitations rendered it unsuitable for sustained ideation work.

\subsubsection{Comparative Observations}

No single model dominated all ideation tasks. Using multiple models proved valuable---each brought distinct strengths to the self-interview process. The practical outcome was a division of labour: GPT-4.5 for initial brainstorming where its occasional insights mattered most, and Claude for sustained structured questioning.

Table~\ref{tab:ideation-comparison} summarises ideation performance across models.

\begin{table}[h]
\centering
\caption{Ideation Performance Comparison}
\label{tab:ideation-comparison}
\begin{tabular}{p{3cm}p{4cm}p{5cm}}
\toprule
\textbf{Model} & \textbf{Strength} & \textbf{Key Limitation} \\
\midrule
GPT-4.5 & Occasional insights & (Synthesis limitations discussed in Section~\ref{sec:lit-discovery}) \\
o1 Pro & (Not used for ideation) & --- \\
o3 & (Limited applicability) & Corporate marketing register; formulaic questions \\
Claude 3.7 Sonnet & Academic register; focused questioning & (No significant ideation-specific limitations observed) \\
Gemini 2.5 Pro & Question quality (vs. o3) & Inconsistent; lacked Claude's focus or GPT-4.5's occasional insights \\
\bottomrule
\end{tabular}
\end{table}

% ----------------------------------------------------------------------
% END BATCH 1
% ----------------------------------------------------------------------

% ======================================================================
% PLACEHOLDER SECTIONS - Batches 2-4
% ======================================================================

% ----------------------------------------------------------------------
% 4.3 LITERATURE DISCOVERY AND EVALUATION (Batch 2)
% ----------------------------------------------------------------------

\subsection{Literature Discovery and Evaluation}
\label{sec:lit-discovery}

We evaluated multiple platforms for systematic literature discovery on software longevity and sustainability in digital humanities, as well as evaluation and synthesis of sources. Systems varied in their ability to surface relevant literature. All systems, furthermore, tended towards uncritical acceptance of source claims rather than interrogation of their validity. This credulity manifested differently across platforms but represented a consistent failure of scholarly judgment.

\subsubsection{OpenAI ChatGPT Deep Research}

ChatGPT Deep Research identified 43 valid sources (37 unique to this system) from a corpus of 93 total unique sources across all methods. ChatGPT required two prompting attempts: initial results using a standard prompt produced unusable outputs, while a revised prompt emphasising scholarly sources yielded the documented results. Of 58 total sources returned, 4 were irrelevant to the research topic and 11 could not be located despite extensive searching and likely represent hallucinations or confaulations.

The service's output composition included 20 journal articles, 13 reports, 4 conference papers, 2 academic blog posts, 2 preprints, 1 presentation, and 1 web page among valid sources. All sources met academic seriousness thresholds, with no lightweight or non-academic content requiring exclusion. Citation accuracy proved problematic: 10 sources had incorrect URLs requiring manual discovery, at least 5 contained incorrect DOIs, and at least 13 featured poorly formatted or incomplete author information. When we requested BibTeX export of discovered sources, ChatGPT only provided it for part of the corpus, despite repeated prompts. ChatGPT could not be convinced to undertake ``citation chaining'' (mining the bibliography of discovered sources for additional sources), a limitation shared with other models.

Five sources overlapped with Perplexity's results, including foundational works on digital humanities sustainability by Barats et al.\ (2020) and Tucker (2022). One source---a systematic literature review on software sustainability by Imran and Kosar (2019)---was also identified by Elicit. The service's 21 sources on national and international initiatives came from targeted follow-up queries, demonstrating responsiveness to prompt refinement.

\subsubsection{Other Tools}

Perplexity Research identified 20 valid sources (14 unique). The initial standardised query produced 9 valid sources alongside 6 invalid items, while follow-up queries yielded the remaining 11 valid sources---demonstrating diminishing returns as subsequent searches produced proportionally more irrelevant results. Of 34 total items returned, 14 were deemed invalid: 5 web pages failed to meet academic seriousness thresholds, 8 sources proved irrelevant to the research topic despite potential value in other contexts, and 1 article appeared in a predatory journal. The service exhibited the highest proportion of non-academic and irrelevant sources among tested platforms, with follow-up queries often losing research context entirely, producing results about oil field decommissioning or end-of-life care rather than software sustainability.

Elicit identified 10 valid sources (8 unique)---the smallest yield among tested platforms. The service returned no invalid or unusable references, demonstrating superior precision despite limited coverage. Citation accuracy exceeded other tested services, though 3 of 10 records contained incorrectly formatted or incomplete author information. The service's reliance on Semantic Scholar limited its discovery scope but contributed to greater precision in source selection.

Claude Research and Gemini Deep Research were briefly tested for literature discovery after the primary literature discovery period (February 2025) but did not surfce additional, useful bibliography. Gemini, for example, returned fewer sources meeting quality thresholds for inclusion, with high confabulation rates; it frequently invented plausible-sounding citations that did not correspond to real publications. Unlike other systems' useful near-misses that could lead to relevant literature, Gemini's offered no serendipitous value. Examination of its ``thinking'' logs revealed a tendency to conflate multiple complex tasks into single operations. In one Digital Humanities search, Gemini declared: ``My next steps involve a more systematic approach to understanding the academic landscape by exploring a directory of DH centers. I also plan to analyze a key chapter discussing the future of DH research.'' It then claimed to have ``successfully analysed a chapter from `Debates in the Digital Humanities''' after downloading only a snippet. Surprisingly, Gemini did not utilise Google Scholar, a missed opportunity. Gemini's results represented what we termed ``source soup''---collecting search results without meaningful evaluation.

\subsubsection{Source Synthesis}

When tasked with synthesising discovered literature into coherent analysis, all tested models exhibited significant limitations. GPT-4.5 defaulted to sequential summarisation rather than analytical integration. The o1 Pro model showed only incremental improvements in instruction adherence. Both models exhibited what we term ``confidence collapse'' at synthesis boundaries, retreating to meta-commentary (``scholars debate this point'') or manufacturing false consensus when confronted with conflicting evidence. Claude 3.7 Sonnet demonstrated better adherence to synthesis instructions but still tended to accept claims at face value, failing to interrogate underlying assumptions or methodological rigor. Gemini 2.5 Pro's synthesis attempts were superficial, often reiterating source content without meaningful integration or critique. Ultimately, models were used to discover sources, but did not contribute significantly to analytical synthesis of sources in our literature review.

\subsubsection{Comparative Performance}

Our investigation revealed marked differences between traditional academic databases and AI-assisted search tools. The complete corpus comprised 93 unique sources addressing software sustainability and longevity in digital humanities and related fields. The three AI services tested for literature discovery collectively contributed 66 sources (71.0\%), while traditional methods contributed 27 (29.0\%). Traditional methods included use of scholarly databases like Web of Science, serendipitous discoveries during AI reference correction, citation chaining, and expert knowledge. Web of Science returned 50 results of which only 3 proved relevant, suggesting fundamental limitations in faceted search when addressing topics that span disciplinary boundaries. Traditional databases do, however, facilitate citation chaining, a technique that none of the AI models could be convinced to apply, even when specifically prompted to do so. 

\begin{table}[h]
\centering
\caption{Comparative Literature Discovery Performance}
\label{tab:literature-performance}
\begin{tabular}{lccccc}
\toprule
\textbf{Method} & \textbf{Total} & \textbf{Unique} & \textbf{\% of Corpus} & \textbf{Required Correction} & \textbf{Errors} \\
\midrule
ChatGPT Deep Research & 58 & 37 & 39.8\% & $\sim$67\% & 15 \\
Perplexity Research & 34 & 14 & 15.1\% & $\sim$67\% & 14 \\
Elicit & 10 & 8 & 8.6\% & $\sim$30\% & 0 \\
Traditional methods & 27 & 27 & 29.0\% & --- & 0 \\
Multi-service & --- & 7 & 7.5\% & --- & --- \\
\midrule
\textbf{Total} & 129 & 93 & 100\% & --- & 29 \\
\bottomrule
\end{tabular}
\smallskip
{\small \textit{Notes}: \textit{Total}: sources added to reference library from this method. \textit{Unique}: sources found only by this method. \textit{\% of Corpus}: proportion of the 93-source corpus contributed uniquely by this method. \textit{Required Correction}: approximate proportion of valid sources needing manual correction of metadata (URLs, DOIs, author formatting). \textit{Errors}: sources deemed invalid (irrelevant, non-academic, predatory, not locatable, or confabulated). Traditional methods include use of scholarly databases like Web of Science, citation chaining, and serendipitous discoveries during reference correction.}
\end{table}

\begin{table}[h]
\centering
\caption{Literature Discovery Error Breakdown by Service and Error Type}
\label{tab:literature-errors}
\begin{tabular}{lccccc}
\toprule
\textbf{Service} & \textbf{Non-academic} & \textbf{Not-relevant} & \textbf{Predatory} & \textbf{Confabulated} & \textbf{Total} \\
\midrule
ChatGPT Deep Research & 0 & 4 & 0 & 11 & 15 \\
Perplexity Research & 5 & 8 & 1 & 0 & 14 \\
Elicit & 0 & 0 & 0 & 0 & 0 \\
\midrule
\textbf{Total} & 5 & 12 & 1 & 11 & 29 \\
\bottomrule
\end{tabular}
\smallskip
{\small \textit{Notes}: \textit{Non-academic}: sources failing to meet scholarly standards (e.g., commercial blogs, marketing content). \textit{Not-relevant}: academically serious sources outside the research scope. \textit{Predatory}: sources from predatory journals. \textit{Confabulated}: citations that could not be located despite extensive searching, likely representing hallucinations or malformed references.}
\end{table}

The various LLMs demonstrated remarkably little overlap in their discoveries. No source appeared in all three services' results. Only seven sources were identified by multiple platforms. This minimal overlap suggests that each service employs distinct search strategies and source prioritisation, potentially offering complementary rather than redundant coverage.

Citation quality emerged as a consistent challenge across all AI platforms. Approximately two-thirds of references from ChatGPT Deep Research and Perplexity required manual correction, with errors ranging from incorrect author formatting and invalid DOIs to erroneous URLs. Even Elicit, which demonstrated superior precision, produced errors in 30\% of its citations. These error rates necessitated manual verification of every AI-generated reference---a time-intensive process that offset much of the efficiency gained through automated discovery.

% ----------------------------------------------------------------------
% 4.4 TOOL DISCOVERY (Batch 3)
% ----------------------------------------------------------------------

\subsection{Tool Discovery}
\label{sec:tool-discovery}

The core of the software longevity study was an empirical examination of research software tools, including discovery of tools, documentation of tools with metadata, and gathering of 'evidence of life' of a tool (e.g., mention of the tool or examples of its use). LLMs assisted with all three stages. 

First, we tasked AI systems with discovering software tools used in archaeology and historical sciences, drawing from academic journals including Internet Archaeology, the Journal of Open Source Software (JOSS), the Journal of Computer Applications in Archaeology (JCAA), the Journal of Open Archaeology Data (JOAD), and SoftwareX. This task tested the systems' ability to discover domain-specific research tools from scholarly literature.

\subsubsection{Discovery Results}

Across multiple discovery runs against the target journals, the systems identified 242 unique tools, of which 154 were verified as legitimate research software. Not all discoveries were valid: 52 were misattributions (real software that did not meet the definition of `tool' provided to the LLM) and 35 were confabulations (fabricated tools with no real-world counterpart). ChatGPT Deep Research served as the primary discovery engine, contributing 122 of the 129 uniquely-discovered verified tools (94.6\%). OpenAI's o3 model was tested on a subset of journals but produced high error rates; Perplexity Deep Research and OpenAI Operator provided limited additional coverage. Gemini and Claude did not yet have dedicated research capabilities during the discovery period and were not tested. Table~\ref{tab:discovery-by-model} summarises outcomes by model.

\begin{table}[h]
\centering
\caption{Tool Discovery by Model}
\label{tab:discovery-by-model}
\begin{tabular}{lcccccc}
\toprule
\textbf{Model} & \textbf{Total} & \textbf{Unique} & \textbf{Verified} & \textbf{Unique Ver.} & \textbf{Misattrib.} & \textbf{Confab.} \\
\midrule
ChatGPT Deep Research & 227 & 188 & 146 & 122 (94.6\%) & 45 & 35 \\
OpenAI o3 & 46 & 12 & 24 & 5 (3.9\%) & 8 & 14 \\
Perplexity Deep Research & 9 & 0 & 9 & 0 & 0 & 0 \\
OpenAI Operator & 3 & 2 & 3 & 2 (1.6\%) & 0 & 0 \\
Multi-model & --- & 40 & 25 & --- & --- & --- \\
\midrule
\textbf{Total} & --- & \textbf{242} & \textbf{154} & \textbf{129} & \textbf{52} & \textbf{35} \\
\bottomrule
\end{tabular}
\smallskip
{\small \textit{Notes}: \textit{Total}: all tools found by this model (includes overlaps). \textit{Unique}: tools found only by this model. \textit{Verified}: legitimate research software confirmed in claimed sources. \textit{Unique Ver.}: unique verified tools contributed by this model (percentage of 129 total). \textit{Misattributed}: real software mentioned in a different context than claimed. \textit{Confabulated}: fabricated tools with no corresponding real software or publication. Multi-model row shows tools found by two or more models.}
\end{table}

Discovery outcomes varied substantially by journal. Internet Archaeology, with its 30-year publication history and extensive web presence, yielded the most verified tools (84 unique) with a confabulation rate below 1\%. JOSS contributed the second-largest yield (48 unique verified tools) but exhibited a 24\% confabulation rate. JOAD proved highly problematic, with an 83\% confabulation rate that rendered most discoveries unusable; confabulation patterns are examined below. JCAA and SoftwareX produced moderate yields with varying error profiles. Table~\ref{tab:discovery-by-journal} summarises outcomes by journal.

\begin{table}[h]
\centering
\caption{Tool Discovery by Journal}
\label{tab:discovery-by-journal}
\begin{tabular}{lcccccc}
\toprule
\textbf{Journal} & \textbf{Total} & \textbf{Unique} & \textbf{Verified} & \textbf{Unique Ver.} & \textbf{Misattrib.} & \textbf{Confab.} \\
\midrule
Internet Archaeology & 130 & 127 & 86 & 84 (54.5\%) & 42 & 1 \\
JOSS & 67 & 67 & 48 & 48 (31.2\%) & 3 & 16 \\
JOAD & 18 & 17 & 3 & 2 (1.3\%) & 0 & 15 \\
JCAA & 17 & 13 & 9 & 6 (3.9\%) & 8 & 0 \\
SoftwareX & 12 & 12 & 9 & 9 (5.8\%) & 0 & 3 \\
Multi-journal & --- & 6 & 5 & 5 (3.2\%) & --- & 0 \\
\midrule
\textbf{Total} & --- & \textbf{242} & \textbf{154} & \textbf{154} & \textbf{52} & \textbf{35} \\
\bottomrule
\end{tabular}
\smallskip
{\small \textit{Notes}: \textit{Total}: all tools found in this journal (includes overlaps). \textit{Unique}: tools found only in this journal. \textit{Verified}: legitimate research software confirmed in claimed sources. \textit{Unique Ver.}: unique verified tools contributed by this journal (percentage of 154 total). \textit{Misattributed}: real software mentioned in a different context than claimed. \textit{Confabulated}: fabricated tools with no corresponding real software or publication. Multi-journal row includes tools found in multiple journals or in sample runs without single-journal attribution.}
\end{table}

Multiple discovery runs were conducted for each journal as prompts were refined and different models tested. Tool counts represent the union of all runs, deduplicated at the tool level; verification outcomes were assigned per-tool rather than per-run, allowing early unsuccessful runs to contribute valid discoveries alongside their errors. Error rates improved substantially across the campaign: early exploratory runs exhibited approximately 52\% error rates; main production runs with structured prompts achieved approximately 33\%; and final secondpass runs with extensive constraints achieved 7.4\%. The implications of this prompt-dependent reliability are discussed in Section~\ref{sec:discussion}.

\subsubsection{Confabulation Patterns}

Two distinct confabulation patterns emerged during verification. The first, ``wholesale fabrication,'' appeared most strikingly in the JOAD run. Of 18 tools purportedly discovered from JOAD, 15 proved to be complete fabrications: the tool names, article titles, authors, and DOIs were all invented. These fabrications followed plausible naming conventions---tools like ``3DArq,'' ``ArtifactML,'' ``CeramicNet,'' and ``DataArq'' sound like legitimate archaeology software---but did not correspond to real tools. The fabricated DOIs followed JOAD's actual DOI pattern (10.5334/joad.xxx) but pointed to non-existent articles.

Examination of the chat transcript revealed the mechanism behind this failure. When prompted to continue from JOAD issues 1--8 (which used Deep Research mode and returned legitimate articles) to issues 9--16, the model silently dropped out of research mode. Its internal reasoning trace explicitly stated: ``I need to simulate a plausible list of articles for JOAD issues 9-16\ldots\ I'll create sample article titles.'' The resulting output included fabricated author names (``Alice Brown; Bob Smith,'' ``Emily Green; Frank Taylor''), invented article titles, and fake tool names---all formatted to match legitimate JOAD output. Only when the researcher noticed the absence of the Deep Research indicator and challenged the response (``Your prior response wasn't deep research'') did the model acknowledge the issue and produce actual search results. This incident demonstrates how confidently-presented fabrication can occur without explicit warning when agentic capabilities silently fail.

The second pattern, which we term ``DOI sequence walking,'' appeared in JOSS discovery runs. Both o3 and ChatGPT Deep Research generated tool names attached to DOIs that followed JOSS's sequential numbering pattern but did not correspond to the claimed software. Tools like ``ChronCluster,'' ``ChronoModelr,'' and ``FQC'' were presented with JOSS DOI patterns (10.21105/joss.xxxx) that, when checked, resolved to entirely different publications. This pattern suggests the models recognised the DOI structure and extrapolated plausible entries rather than verifying actual journal contents.

A troubling finding emerged when both o3 and ChatGPT Deep Research independently generated identical fabricated tool names for JOSS. Rather than indicating cross-validation of legitimate discoveries, this convergence revealed shared confabulation modes---both models apparently drew on similar training data or reasoning patterns that produced the same false positives. This finding undermines any assumption that agreement between models indicates reliability.

\subsubsection{Verification Outcomes}

Manual verification of all 242 discovered tools revealed the following distribution: 154 tools (64\%) were verified as legitimate research software mentioned in the claimed sources; 52 tools (21\%) were misattributed, representing real software mentioned in a different context than claimed; and 35 tools (14\%) were confirmed confabulations with no corresponding real software or publication. One tool was classified as a granularity error (a generic software category reported as a specific tool).

The stark variation in confabulation rates across sources---from below 1\% for Internet Archaeology and 0\% for JCAA to 83\% for JOAD---suggests that model reliability depends heavily on the specific corpus being queried. Journals with extensive web presence and consistent metadata structures (Internet Archaeology) yielded reliable results, while those with smaller publication counts or less web visibility (JOAD) triggered wholesale fabrication. This variability means researchers cannot assume consistent performance even within a single discovery task.

% ----------------------------------------------------------------------
% 4.5 TOOL DOCUMENTATION / METADATA (Batch 4)
% ----------------------------------------------------------------------

\subsection{Tool Documentation}

% TODO: Write fresh from quantitative data
% Key findings: 98 tools investigated, 90 successful (92%), 8 failed

% SAVED FROM 4.2 - Claude's scope assessment (relevant to metadata task):
% A revealing pattern emerged in Claude's scope assessment. Unlike Deep Research's
% nine-stage plans, Claude made realistic claims about processing capacity, stating
% ``if you want something thorough, I can do 1 to 3'' tools and explicitly warning
% that ``with 10 tools, you're going to get a shallow summary.'' This operational
% self-awareness did not extend to content evaluation, where Claude remained as
% credulous as its competitors regarding source claims.

% ----------------------------------------------------------------------
% 4.6 TOOL EVIDENCE COLLECTION (Batch 4)
% ----------------------------------------------------------------------

\subsection{Tool Evidence Collection}

% TODO: Write fresh from quantitative data
% Key findings: 1,179 evidence rows, verification outcomes (91% confirmed, 5% confabulated)

